\chapter*{Conclusioni}
\markboth{Conclusioni}{}

In questo lavoro è stato progettato, implementato e validato un framework automatico per la compressione e l'ottimizzazione di modelli \textit{Vision Transformer} (ViT), con l'obiettivo di ridurne l'impatto computazionale e renderne possibile l'esecuzione efficiente su architetture hardware a risorse limitate, come i dispositivi mobili ed \textit{edge}. La strategia proposta ha integrato una tecnica di \textit{pruning} strutturato con gli algoritmi di ottimizzazione, unendo anche la \textit{Knowledge Distillation}. \par

L'analisi sperimentale, condotta su benchmark standard quali \textit{CIFAR-100} e \textit{ImageNet-1K}, ha confermato l'efficacia dell'approccio e la sua competitività con lo stato dell'arte. I risultati quantitativi hanno evidenziato che sia possibile ottenere una riduzione di circa un terzo dei parametri e dei GFLOPs, con una perdita di accuratezza trascurabile, che corrisponde però ad un incremento significativo del throughput. Sono stati inoltre condotti dei test di inferenza su dispositivi mobili, dimostrando l'impatto reale della compressione sulle latenze registrato a seguito della procedura di \textit{deploy}. \par

Gli studi di ablazione hanno validato l'architettura interna del framework, dimostrando che sia la metrica di importanza basata sul Fisher Score, sia la componente di ricerca contribuiscono ad evitare il decadimento delle prestazioni durante il processo di compressione. \par

Nell'analisi dei risultati ottenuti dai test sono emersi alcuni pattern strutturali interessanti che portano ad alcune considerazioni sulla natura dei modelli ViT e sul loro funzionamento. Da quanto osservato emerge una struttura ad "U", già teorizzata e studiata nella letteratura recente, secondo il quale i layer centrali contengono le rappresentazioni ad alto livello, i concetti semantici, e di conseguenza sono più sensibili al pruning. 


\section{Limitazioni e sviluppi futuri}
Di seguito sono elencate alcune limitazioni del lavoro svolto e possibili sviluppi futuri che potrebbero essere intrapresi per migliorare ulteriormente il framework di compressione proposto:

\begin{itemize}
    \item \textbf{Costo computazionale della compressione:} Il framework di compressione presentato, nonostante l'efficacia dimostrata, richiede un costo computazionale considerevole. Se in particolare si vuole recuperare la massima \textit{accuracy} possibile, è necessario utilizzare un numero elevato di epoche di \textit{fine tuning}, che, se unito ad un elevato numero di iterazioni di ricerca, può portare a costi paragonabili a quelli di un addestramento completo del modello originale. Nonostante ciò, è importante sottolineare che il processo di compressione rappresenta un investimento da eseguire in una singola occasione, che presenta benefici tangibili che potrebbero comunque risultare vantaggiosi in scenari di produzione.
    
    \item \textbf{Mancanza di test sui modelli di grandi dimensioni:} A causa dell'elevato costo computazionale prima esposto, i test sperimentali sono stati condotti su modelli di dimensioni ridotte, che rappresentano comunque un'ottima sfida per la compressione, tuttavia non è stato possibile estendere l'analisi a modelli di maggiori dimensioni, che potrebbero beneficiare maggiormente di un processo di questo tipo. Sarebbe inoltre auspicabile estendere l'analisi a modelli di maggiori dimensioni al fine di verificare la presenza degli stessi pattern strutturali osservati nel test condotti in questo lavoro, in maniera tale da confermare ulteriormente le considerazioni formulate attraverso l'analisi dei risultati.

    \item \textbf{Estensione a compiti di visione complessi:} Oltre a testare il framework su modelli di maggiori dimensioni, sarebbe utile eseguire dei test al di fuori del \textit{task} di classificazione. Questo consentirebbe di verificare la bontà del processo su modelli utilizzati come \textit{backbone} per compiti di \textit{object detection}, \textit{semantic segmentation} o modelli multimodali.
    
    \item \textbf{Confronto con metodo Greedy:} Dato l'impiego di un algoritmo di ricerca avanzato per l'individuazione delle migliori azioni di \textit{pruning} da attuare, sarebbe necessario un confronto con un implementazione di tipo \textit{greedy} del metodo. Tale comparazione permetterebbe di verificare una eventuale superiorità del processo di ricerca con algoritmo \textit{Branch and Bound} rispetto a quanto proposto in questo lavoro, giustificando così l'incremento del costo computazionale.

    \item \textbf{Integrazione con tecniche di quantizzazione:} Il beneficio computazionale potrebbe essere ulteriormente amplificato combinando il \textit{pruning} strutturato con tecniche di quantizzazione post-addestramento (PTQ) o quantizzazione di tipo \textit{training-aware}(QAT). L'integrazione di queste tecniche, oggi molto diffuse nell'ambito dei \textit{Large Language Models} (LLM) e pienamente supportate dall'hardware, permetterebbero di ottenere modelli ancora più performanti, al costo di un leggero decadimento delle prestazioni predittive.
    
    \item \textbf{Applicazione a modelli di tipo \textit{Large Language Models} (LLM):} Un interessante sviluppo futuro potrebbe essere rappresentato dall'estensione del framework a modelli di tipo LLM, che presentano un numero di parametri estremamente elevato, e che potrebbero beneficiare in maniera significativa da un processo di compressione. L'applicazione del framework a questo tipo di modelli potrebbe richiedere alcune modifiche alla strategia di compressione, ma rappresenterebbe un importante passo avanti. La più grande limitazione in questo caso sarebbe costituita dal costo computazionale, che potrebbe essere proibitivo, tuttavia si potrebbe mitigarlo attraverso l'impiego di tecniche di \textit{parameter efficient fine tuning} (PEFT), che permetterebbero di ridurre significativamente il numero di parametri da aggiornare durante la fase di \textit{fine tuning}.
    
    \item \textbf{Utilizzo di tecniche di Token Pruning:} Un ulteriore direzione di ricerca potrebbe essere rappresentato dall'integrazione di tecniche di \textit{token pruning}, che permetterebbero di ridurre in maniera significativa il costo computazionale durante la fase di inferenza, eliminando i token meno rilevanti per la predizione. Questo è facilmente visibile nei risultati ottenuti in \cite{wang2025caittriplewincompressionhigh}, nel quale è stato adottato sia un processo di \textit{pruning} strutturato, sia una tecnica di \textit{token merging}, ottenendo un incremento del 90\% di \textit{throughput}.
    
\end{itemize}


